\Needspace{6\baselineskip}
\begin{center}\hypertarget{native-258-article185}{}\pdfbookmark[2]{185.}{article-185}185.\end{center}
The subject we have dealt with so far could give rise to several elegant analytical investigations. We shall refrain from dwelling on them here for fear of being drawn too far from our purpose. For the same reason we must reserve for another occasion an account of the techniques by which the numerical calculation can be reduced to a more efficient algorithm. We may add one observation here. Whenever the number of \SourceLeaf{259}given functions or equations is considerable, the calculation becomes somewhat more troublesome, chiefly because the coefficients by which the original equations must be multiplied to obtain $P$, $Q$, $R$, $S$, etc. usually involve unwieldy decimal fractions. If in this case it does not seem worthwhile to carry out these multiplications as accurately as possible with the aid of logarithmic tables, in most cases it will suffice to replace these multipliers by others that are more convenient for calculation and differ little from them. This latitude can produce appreciable errors only when the measure of precision in the determination of the unknowns turns out to be much smaller than the precision of the original observations.
